\UseRawInputEncoding
%% -*- coding: utf-8 -*-
\documentclass[12pt,a4paper]{scrartcl} 
\usepackage[utf8]{inputenc}
\usepackage[T2A]{fontenc} 
\usepackage[english,russian]{babel}
\usepackage{indentfirst}
\usepackage{misccorr}
\usepackage{graphicx}
\usepackage{amsmath}
\usepackage{listings} 
\lstset{
  basicstyle=\ttfamily\small,
  breaklines=true,
  showstringspaces=false,
  extendedchars=true
}
\begin{document}
\section{Введение}
\label{sec:intro}
\begin{enumerate}
 \item \textbf{Текстовая формулировка задачи:} Нахождение ранга матрицы методом Гаусса.
 \item \textbf{Пример кода:} Представлен в разделе~\ref{sec:exp:code}.
 \item \textbf{График:} Для данного метода используется матричное представление.
 \item \textbf{Скриншот программы:} Представлен в разделе~\ref{sec:picexample}.
\end{enumerate}
\section{Ход работы}
\label{sec:exp}

\subsection{Код приложения}
\label{sec:exp:code}
\begin{lstlisting}[language=Python]
import numpy as np
def get_matrix_from_input():
    try:
        rows = int(input("Enter the number of lines: "))
        if rows <= 0: raise ValueError("The number of rows should be > 0")
        matrix = []
        for i in range(rows):
            row = list(map(float, input(f"Enter the elements of the {i+1}-th row separated by spaces: ").split()))
            matrix.append(row)
            
        if len(set(len(r) for r in matrix)) > 1:
            raise ValueError("All lines must have the same length!")
        return matrix
    except ValueError as e:
        print(f"Input error: {e}")
        return None
def calculate_rank(matrix):
    A = np.array(matrix, dtype=float)
    rows, cols = A.shape
    rank = 0
    pivot_row = 0
    for j in range(cols):
        if pivot_row >= rows:
            break
            
        max_idx = np.argmax(np.abs(A[pivot_row:, j])) + pivot_row
        if np.isclose(A[max_idx, j], 0):
            continue
        A[[pivot_row, max_idx]] = A[[max_idx, pivot_row]]
        
        for i in range(pivot_row + 1, rows):
            factor = A[i, j] / A[pivot_row, j]
            A[i, j:] -= factor * A[pivot_row, j:]
            
        pivot_row += 1
        rank += 1
    return rank
matrix_data = get_matrix_from_input()
if matrix_data is not None:
    rank = calculate_rank(matrix_data)
    print(f"The rank of the matrix is: {rank}")
\end{lstlisting}
\subsection{Математическая постановка}
\label{sec:mathexample}
Ранг матрицы определяется через приведение к трапецеидальному виду:
\begin{equation}\label{eq:matrix}
 A_{m \times n} = \begin{pmatrix} 
 a_{11} & a_{12} & \dots & a_{1n} \\
 0 & a_{22} & \dots & a_{2n} \\
 \dots & \dots & \dots & \dots \\
 0 & 0 & \dots & 0
 \end{pmatrix}
\end{equation}
\section{Пример вставки изображения}
\label{sec:picexample}
\begin{figure}[h]
    \centering
    \includegraphics[width=0.5\textwidth]{BeginCode.png} 
    \caption{Трапециевидная матрица}\label{fig:matrix}
\end{figure}

\section{Список литературы}
\begin{thebibliography}{9}
\bibitem{Knuth-2003}Кнут Д.Э. Всё про \TeX. --- Москва: Изд. Вильямс, 2003 г. [cite: 4]
\bibitem{Lvovsky-2003}Львовский С.М. Набор и верстка в системе \LaTeX{}. --- 3-е издание, 2003 г. [cite: 4]
\bibitem{Voroncov-2005}Воронцов К.В. \LaTeX{} в примерах. 2005 г. [cite: 4]
\end{thebibliography}
\end{document}